% Modernised reading — fonds Alexandre Grothendieck, shelfmark 29, batch 2
% (pages 21-40).
%
% This is an INTERPRETATION, not a transcription. It re-reads the pages in
% today's notation and vocabulary, states the hypotheses the manuscript leaves
% implicit, and drops the critical apparatus so that the mathematics reads
% straight through. Everything the brackets and underlines carried moves into a
% footnote. In any dispute of substance the transcription governs: whoever
% cites Grothendieck must cite batch-02.fr.tex.
%
% Scope: the folder was read WHOLE before any of this was written — all eleven
% batches, pages 1-216, in one pass — and one file was then written per batch.
% Two things in this file are only visible at that scope: the missing
% hypothesis in the definition of a Galois covering of type R, supplied here
% from page 36; and the word « multigaloisienne », used on page 30 as if
% already available and axiomatised only on pages 208-216.
%
% Notation is the folder's, fixed in batch 1's reading: Et(S), Fl_S,
% calligraphic R = (C, s) for the 1967 *donnée*, Rev^R(S) with the exponent,
% r for the geometric realisation.
%
% Departures from the page, each footnoted where it occurs:
%   - page 30's definition of a Galois covering of type R is completed by the
%     commutation condition Grothendieck states only on page 36;
%   - page 21 writes « sur B » for « sur S »;
%   - the profinite group of page 25 requires the category to be GALOISIAN,
%     i.e. the connectedness he imposes; without it one gets a groupoid, which
%     is what « multigaloisienne » on page 30 means.
%
% Pass: Opus 5 (claude-opus-5), 2026-08-29 — first pass, unchecked against the
% pages by a human.
\documentclass[11pt,a4paper]{article}
% Le préambule commun à toutes les transcriptions du fonds Grothendieck.
%
% Il tient en une idée : **l'incertitude se compose, elle ne se résout pas.**
% Une transcription qui lisse un mot douteux a détruit la seule chose pour
% laquelle on la faisait — un lecteur doit pouvoir savoir, à chaque endroit,
% ce qui était sur le feuillet et ce qui a été ajouté. Les six macros qui
% suivent sont donc l'appareil critique, et non de la décoration.
%
% Les mêmes commandes sont reconnues par `scripts/render.mjs`, qui produit la
% vue de lecture du site. Une macro ajoutée ici doit l'être aussi là-bas,
% faute de quoi le rendu échouera bruyamment — ce qui est voulu : un rendu
% silencieusement faux est pire qu'un rendu absent.

\ProvidesPackage{grothendieck}[2026/08/08 Transcriptions du fonds Grothendieck]

\usepackage{amsmath,amssymb,amsthm}
\usepackage{mathrsfs}
\usepackage{stmaryrd}
% Les diagrammes commutatifs sont partout dans ce fonds : c'est la langue
% même de la géométrie algébrique de Grothendieck, et une transcription qui
% les rendrait en prose ne transcrirait rien.
\usepackage{tikz-cd}
% Français par défaut — c'est la langue des feuillets — mais l'anglais est
% chargé aussi : la traduction et le résumé sont des documents anglais, et la
% typographie française (espaces avant les deux-points) y serait une faute.
% Ces documents basculent par \selectlanguage{english} en tête de corps.
\usepackage[english,french]{babel}
\usepackage{fontspec}
\usepackage{geometry}
\geometry{a4paper,margin=25mm,marginparwidth=28mm}
\usepackage{marginnote}
\usepackage{xcolor}
% ulem, not soul. `soul`'s \st and \ul are built on a character-by-character
% scan that XeLaTeX's font machinery breaks: the document fails with an
% undefined control sequence rather than degrading. ulem does the same two jobs
% and is Unicode-engine safe. `normalem` stops it hijacking \emph, which the
% transcriptions use for Grothendieck's own emphasis.
\usepackage[normalem]{ulem}
\usepackage{hyperref}
\hypersetup{colorlinks=true,linkcolor=black,urlcolor=black,pdfborder={0 0 0}}

% --- Métadonnées du lot -----------------------------------------------------
% Reprises telles quelles par le script de rendu, qui les affiche en tête de
% la vue de lecture. Elles fixent ce que le fichier prétend couvrir : sans
% elles, une transcription flotte sans point d'attache dans un fonds de seize
% mille feuillets.

\newcommand{\folder}[1]{\gdef\grothFolder{#1}}
\newcommand{\batch}[1]{\gdef\grothBatch{#1}}
\newcommand{\leaves}[2]{\gdef\grothFirst{#1}\gdef\grothLast{#2}}
\newcommand{\foldertitle}[1]{\gdef\grothFoldertitle{#1}}
\gdef\grothFolder{?}\gdef\grothBatch{?}\gdef\grothFirst{?}\gdef\grothLast{?}\gdef\grothFoldertitle{}

% --- Le résumé --------------------------------------------------------------
%
% Il ouvre la lecture modernisée, sous le titre « Résumé », et il est composé
% en petit. Ce n'est pas un ornement : c'est le seul passage du document qui
% ne soit pas des mathématiques, et un lecteur doit pouvoir le reconnaître
% avant de l'avoir lu — puis le sauter s'il connaît déjà le sujet, ou s'y
% arrêter s'il ne le connaît pas. Le filet qui le suit marque la reprise du
% corps.

\newenvironment{resume}
  {\small\setlength{\parskip}{0.45em}}
  {\par\medskip\hrule\medskip}

% --- Mots-clés ---------------------------------------------------------------
%
% En anglais, en clôture du résumé de la lecture modernisée : le vocabulaire
% moderne sous lequel chercher ce que les feuillets construisent. Le manifeste
% du site (`npm run manifest`) les extrait pour étiqueter la cote entière —
% cette ligne est la source unique des tags, il n'y a pas de fichier à côté.
\newcommand{\keywords}[1]{\par\smallskip\noindent
  {\footnotesize\textsc{Keywords} --- \textit{#1}}\par}

% --- L'appareil critique ----------------------------------------------------

% Le numéro de feuillet, en marge. C'est l'ancre qui permet de régler un
% désaccord sur une formule en regardant *un* feuillet plutôt qu'une chemise
% de six cents. Le site s'en sert aussi pour tourner les pages du fac-similé
% quand on fait défiler la transcription.
\newcommand{\leaf}[1]{%
  \marginnote{\footnotesize\color{gray!70}\textsf{#1}}%
  \label{leaf:#1}\ignorespaces}

% Illisible : on ne devine pas. Un mot inventé qui se lit comme les autres est
% le pire résultat possible.
\newcommand{\ill}{\textcolor{red!60!black}{[\,\ldots\,]}}

% Lecture douteuse : le mot est donné, et signalé comme incertain.
\newcommand{\uncertain}[1]{\uline{#1}}

% Ajout de l'éditeur — une lettre manquante, un mot sous-entendu.
\newcommand{\add}[1]{\textcolor{blue!50!black}{[#1]}}

% Biffé par Grothendieck lui-même. Ce qu'il a rayé fait partie du document :
% une rature dit souvent où la pensée a changé de direction.
\newcommand{\struck}[1]{\sout{#1}}

% Note du transcripteur, sur l'état matériel du feuillet.
\newcommand{\note}[1]{\par\smallskip{\small\color{gray!60!black}\textit{[#1]}}\par\smallskip}

% Note marginale de Grothendieck, distincte du corps du texte.
\newcommand{\marginal}[1]{\par\smallskip\noindent\hspace{1em}%
  {\small\color{gray!60!black}#1}\par\smallskip}

% --- Titre ------------------------------------------------------------------

\AtBeginDocument{%
  \begin{center}
    {\large\bfseries Fonds Alexandre Grothendieck --- cote n\textsuperscript{o}~\grothFolder}\\[2pt]
    {\itshape\grothFoldertitle}\\[4pt]
    {\small feuillets \grothFirst--\grothLast\ (lot \grothBatch)}\\[2pt]
    {\footnotesize\color{gray!60!black}Transcription établie d'après le fac-similé numérisé
     par l'Université de Montpellier.\\
     Les lectures incertaines sont soulignées, les illisibles notées [\,\ldots\,],
     les ajouts entre crochets.}
  \end{center}
  \vspace{1em}}

\endinput


\folder{29}
\batch{2}
\pages{21}{40}
\dating{1959-1969}
\watermark{Édition de démonstration\\interprétation personnelle de l'œuvre}
\foldertitle{Groupe fondamental [Autour de SGA 4 et SGA 7] — lecture modernisée}

\begin{document}

\section*{Résumé}

\begin{resume}
Le lot précédent a monté un appareil : une \emph{donnée de ramification} sur un
schéma $S$, c'est-à-dire un catalogue de modèles de ramification autorisés, et
la catégorie $\mathrm{Rev}^{\mathcal{R}}(S)$ des revêtements qui ne sont pas
plus ramifiés qu'eux. Ces pages achèvent l'appareil, en extraient un groupe, et
disent --- c'est le plus important --- ce qu'il faut en penser.

Ce qu'il faut en penser est ceci : \emph{la donnée de ramification n'est qu'un
échafaudage}. Grothendieck l'écrit page 22 en toutes lettres. L'objet
véritable, c'est le foncteur qui, à un revêtement de type $\mathcal{R}$,
associe le revêtement ordinaire qu'il réalise ; deux données très différentes
peuvent définir le même foncteur, et alors elles sont « essentiellement
équivalentes ». Il donne l'analogie juste, et c'est la sienne : entre une
donnée de ramification et la catégorie qu'elle définit, il y a le même rapport
qu'entre un \emph{site} et le \emph{topos} associé. Le site est un choix
commode, le topos est l'objet. Puis il ajoute la phrase que le dossier entier
mettra dix ans à honorer : si l'on veut normaliser la notion pour éliminer
l'arbitraire, « il semble qu'on soit infailliblement conduit à prendre comme
catégorie $\mathcal{C}$ la catégorie $\mathrm{Rev}^{\mathcal{R}}$ elle-même
qu'on se propose de construire ». Les dernières pages du carton feront
exactement cela.

Entre-temps, la machine fonctionne. La catégorie hérite de toutes les
propriétés d'exactitude des revêtements étales ; un point géométrique
convenablement choisi --- hors de « l'ensemble de ramification » --- fournit un
foncteur fibre exact ; et la théorie de Galois abstraite rend un groupe
profini, le $\pi_{1}^{\mathcal{R}}(S,s)$, dont les ensembles finis à opérateurs
sont les revêtements. La fonctorialité en $S$ et en $\mathcal{R}$ est traitée,
avec un aveu de fatigue devant la transitivité --- « jetant provisoirement un
voile pudique sur ces abîmes » --- et deux cas particuliers utiles : le
changement de base étale fini, qui donne un sous-groupe, et le changement de
base entier radiciel surjectif, qui ne change rien.

Puis vient le paragraphe 4, sur les revêtements galoisiens, et avec lui le
passage qui justifie le dossier. Grothendieck l'écrit en anglais, en plein
milieu du tapuscrit français, parce qu'il est en train de répondre à Murre :
s'il a si longtemps refusé de définir les revêtements modérément ramifiés en
dehors du cas galoisien, c'est qu'il n'avait pas vu qu'un tel revêtement n'est
pas un revêtement muni d'une \emph{condition} mais un revêtement muni d'une
\emph{structure} ; dans le cas galoisien, la structure est unique si elle
existe, et l'on peut ne pas la voir. Il ajoute deux avertissements qui valent
d'être médités : tous les morphismes naïfs entre revêtements modérés ne sont
pas admissibles, seuls ceux qui viennent de la catégorie le sont ; et le
produit de deux revêtements modérés n'est \emph{pas} le produit fibré
au-dessus de $S$. Conclusion : « the consideration of the category
$\mathrm{Rev}^{\mathcal{R}}(S)$ is undispensable for a good comprehension of
what is really going on ».

Les deux dernières pages ouvrent un tout autre dossier, celui des théorèmes de
finitude, que les lots 3 et 4 poursuivront. Les noms modernes à chercher sont
\emph{groupe fondamental modéré}, \emph{théorème d'Abhyankar}, \emph{descente
fidèlement plate}, et --- pour l'objet vers lequel tout cela tend ---
\emph{catégorie galoisienne}.

\keywords{fibre functor, tamely ramified covering, normal crossings divisor, Abhyankar's lemma, faithfully flat descent}
\end{resume}

\pagerange{21}{22}
\section*{Le cas kummérien local, et la ramification modérée au sens large (pages 21 et 22)}

La page 20 s'arrêtait au milieu du second cas kummérien. Le voici entier.

Soit $D$ un diviseur sur $S$. On considère le faisceau, sur $\mathrm{Et}(S)$,
associé au préfaisceau qui à $S'$ associe l'ensemble des familles $(D_{i})$ de
diviseurs sur $S'$ telles que $D_{S'} = \sum_{i} D_{i}$ et que, pour tout
$s' \in S'$, les restrictions à $\operatorname{Spec}(\mathcal{O}_{S',s'})$ des
$D_{i}$ non nulles soient deux à deux disjointes. Une section de ce faisceau
donne, localement pour la topologie étale, une situation comme au cas
précédent, ce qui suffit à définir une donnée de ramification sur
$S$.\footnote{La page écrit « une donnée de ramification $\mathcal{R}$ sur
$B$ », pour « sur $S$ » ; c'est une coquille de frappe, corrigée ici. Le
tapuscrit appelle d'ailleurs « canularesque » la condition locale sur les
$(D_{i})$ --- elle n'est là que pour qu'on ait bien un préfaisceau.}

Deux cas particuliers méritent un nom, car $\mathcal{C}$ s'y décrit en termes
du seul $D$ :

\begin{enumerate}
  \item[1\textsuperscript{o})] $S$ localement noethérien, $D$ à croisements
    normaux ; on prend pour $D_{i}$ les composantes irréductibles locales
    régulières de $D$, qui ne sont définies que localement pour la topologie
    étale ;
  \item[2\textsuperscript{o})] $S$ au-dessus d'un préschéma $T$, $D$ à
    croisements normaux \emph{relativement à $T$} ; on prend pour $D_{i}$ les
    composantes locales lisses sur $T$.
\end{enumerate}

Dans le cas 1\textsuperscript{o}, $\mathrm{Rev}^{\mathcal{R}}(S)$ est
équivalente à une sous-catégorie pleine de la catégorie des revêtements de $S$,
et l'on appelle ses objets les \emph{revêtements modérément ramifiés au sens
large} relativement à $D$. C'est une notion strictement plus large que celle de
revêtement modérément ramifié par rapport à la famille réduite au seul $D$, et
c'est le \emph{théorème d'Abhyankar} qui la caractérise : ce sont les
revêtements étales sur $S - \operatorname{supp} D$, normaux aux points
au-dessus de $\operatorname{supp} D$, dont toute composante irréductible en
domine une de $S$, et qui sont modérément ramifiés au sens classique aux points
maximaux de $D$.

C'est le seul endroit du carton où la \emph{structure} redevient une
\emph{propriété}, et il vaut d'être noté comme tel : la caractérisation
d'Abhyankar est ce qui rend le cas des croisements normaux exceptionnellement
maniable, et son absence ailleurs est la raison d'être de tout l'appareil.

Le cas 2\textsuperscript{o} attend encore une vérification --- « sauf erreur
(il faudrait le vérifier) » ---, et l'on parlerait alors de revêtements
« modérément ramifiés pour $D$, relativement à la base $T$ ». Lorsque $T$ est
régulier les deux notions coïncident. L'intérêt du cas relatif est indiqué en
une phrase : c'est avec la ramification modérée au sens large que le théorème
de spécialisation pour le groupe fondamental modéré se démontre par la
technique de SGA 1, pour $S$ propre et lisse sur $T$ et $D$ à croisements
normaux relativement à $T$.\footnote{Ce théorème de spécialisation est celui de
Michèle Raynaud, SGA 1 XIII ; le tapuscrit y renvoie explicitement page 114 du
carton, « exposé de Mme Raynaud ».}

\pagerange{22}{24}
\section*{La donnée est auxiliaire : l'analogie du site et du topos (pages 22 à 24)}

Pour $S'$ variable dans $\mathrm{Et}(S)$ on définit
$\mathrm{Rev}^{\mathcal{R}}(S')$ en restreignant $\mathcal{R}$ à
$\mathrm{Et}(S') = \mathrm{Et}(S)_{/S'}$, d'où une catégorie fibrée
$\underline{\mathrm{Rev}}^{\mathcal{R}}$ sur $\mathrm{Et}(S)$ --- qui est même
un \emph{champ} --- et un foncteur cartésien
\[
  (\ast) \qquad
  r : \underline{\mathrm{Rev}}^{\mathcal{R}} \longrightarrow
  \underline{\mathrm{Rev}} ,
\]
où $\underline{\mathrm{Rev}}$ désigne la catégorie fibrée des revêtements des
objets de $\mathrm{Et}(S)$.\footnote{On suppose ici, pour tout $S'$, que la
réalisation géométrique sur $\mathcal{C}_{S'}$ est à valeurs dans les
revêtements de $S'$, et que le foncteur $r$ sur
$\mathrm{Rev}^{\mathcal{R}}(S')$ est bien défini --- c'est l'hypothèse (H2) du
lot 1, prise fibre à fibre.}

C'est ici que le tapuscrit dit ce qu'il faut penser de tout ce qui précède, et
la phrase doit être citée :

\begin{quote}
Il faut insister sur le fait que pour nous, la « donnée de ramification »
$\mathcal{R}$ est un objet de nature essentiellement auxiliaire, destiné à
définir la catégorie fibrée $\underline{\mathrm{Rev}}^{\mathcal{R}}$ (qui est
même un champ) et le foncteur cartésien $r$ de $(\ast)$ ; c'est $(\ast)$ qui
est l'objet véritable de l'étude.
\end{quote}

On dira donc que deux données $\mathcal{R}$, $\mathcal{R}'$ sont
\emph{essentiellement équivalentes} si elles définissent des foncteurs $(\ast)$
équivalents, par une équivalence induisant l'identité sur
$\underline{\mathrm{Rev}}$.

\subsection*{Trois exemples d'équivalence}

Le premier est formel. Sous les conditions de l'exemple 3, soit $\mathcal{C}'$
une sous-catégorie fibrée de $\mathcal{C}$ telle que tout objet $M$ d'un
$\mathcal{C}_{S'}$ soit localement majoré par un objet de $\mathcal{C}'$, et
munissons $\mathcal{R}'$ de la réalisation géométrique induite : alors
$\mathcal{R}$ et $\mathcal{R}'$ sont équivalentes. Autrement dit,
\emph{seul le crible compte, pas le catalogue}.

Le deuxième compare les exemples 2 et 3 dans le cas kummérien. Si les $D_{i}$
admettent des équations globales $c_{i}$, on peut considérer le système
projectif des $Z^{\underline{n}}_{\underline{c}}$ pour $\underline{n}$
variable, premier à la caractéristique de $S$ (supposée exister), avec le
système projectif des groupes $\mathbb{Z}/n_{i}\mathbb{Z}$ ; si l'on dispose
d'un isomorphisme entre ce dernier et le système des $\mu_{\underline{n}}$ ---
il en existe lorsque $S$ est un schéma sur un corps algébriquement clos --- on
est dans les conditions de l'exemple 2, et la donnée obtenue est équivalente à
la donnée kummérienne de l'exemple 3. Les deux exemples décrivent donc, dans ce
cas, la même chose ; ce qui les distingue est l'existence globale des racines
de l'unité.

Le troisième est une torsion. Soit $\mathcal{R}$ la donnée de l'exemple 1,
définie par $(Z,G)$. Donnons-nous un revêtement principal $P$ de $S$ de groupe
$G'$ opérant à droite, et supposons que $G$ opère à gauche sur $P$, librement,
en commutant à $G'$. Alors $Z' = Z \times^{G} P$, muni de $G'$, définit une
donnée $\mathcal{R}'$ équivalente à $\mathcal{R}$ : on le voit d'abord quand
$P = G'_{S}$ avec l'opération régulière droite et $G \hookrightarrow G'$, puis
par localisation étale.

\subsection*{« On soit infailliblement conduit \dots »}

Et vient alors la phrase qui donne au dossier son programme :

\begin{quote}
On peut se proposer de normaliser la notion de donnée de ramification, pour
éliminer l'arbitraire dans le choix d'un $\mathcal{R}$. Mais alors il semble
qu'on soit infailliblement conduit à prendre comme catégorie $\mathcal{C}$ la
catégorie $\mathrm{Rev}^{\mathcal{R}}$ elle-même qu'on se propose de construire
et d'étudier ! (Comparer avec la situation analogue pour les relations entre
sites et topos, l'objet auxiliaire étant le site, l'objet intrinsèque de nature
géométrique étant le topos associé à ce site.)
\end{quote}

Il faut lire cela comme un problème ouvert posé à soi-même, et il est résolu
dans le même carton, deux cents pages plus loin : le \emph{domaine de
ramification} des pages 190 à 200 est un champ de catégories multigaloisiennes
au-dessus de $\mathrm{Et}(S)$, muni d'une réalisation géométrique continue,
c'est-à-dire exactement la catégorie elle-même prise comme donnée. Le carton ne
date pas ces pages-là et cette lecture ne le fera pas ; mais la boucle est
close, et elle l'est dans un seul carton.

\pagerange{24}{25}
\section*{La catégorie et son groupe fondamental (pages 24 et 25)}

\subsection*{Exactitude}

La catégorie $\mathrm{Rev}^{\mathcal{R}}(S)$ hérite de toutes les propriétés
d'exactitude établies pour les revêtements étales d'un schéma. De façon
précise : les limites projectives finies et les limites inductives finies y
existent et satisfont les mêmes propriétés d'exactitude que dans la catégorie
des ensembles ; tout morphisme s'y factorise en un épimorphisme suivi d'un
monomorphisme ; épimorphismes et monomorphismes sont effectifs ; et tout
monomorphisme est un isomorphisme sur un facteur direct.

Ce dernier énoncé est exactement l'axiome que la page 209 numérotera
\textbf{G3}, et que SGA 1 V numérote de même : \emph{tout morphisme se
factorise en un épimorphisme suivi d'un isomorphisme sur un sommande direct}.
Il est ici obtenu comme une propriété, et il sera là-bas posé comme un axiome ;
c'est le même énoncé, à onze ans et cent quatre-vingts pages de distance.

Un avertissement, qui n'est pas une remarque de détail : la réalisation
géométrique $r$ commute aux limites inductives finies --- en particulier aux
sommes --- mais \emph{pas} aux limites projectives finies, et en particulier
pas au produit de deux objets. Le produit de deux revêtements de type
$\mathcal{R}$ n'est donc pas le produit fibré de leurs réalisations au-dessus
de $S$. C'est le fait que la page 31 signalera comme la cause de la
complication de l'un des lemmes de Murre.

\subsection*{Le foncteur fibre}

Pour appliquer la théorie de Galois abstraite, il faut de plus un
\emph{foncteur fibre}
\[
  F : \mathrm{Rev}^{\mathcal{R}}(S) \longrightarrow (\text{ensembles finis}) ,
\]
c'est-à-dire un foncteur exact.

Dans l'exemple 1, tout point géométrique de $Z$ en fournit un ; il est
\emph{conservatif} si toute partie ouverte et fermée de $Z$ stable par $G$ est
vide ou égale à $Z$. Dans l'exemple 2, on prend un point géométrique de
$Z = \varprojlim Z_{i}$.

Dans le cas général de l'exemple 3, il faut travailler un peu. On choisit un
point géométrique $s$ de $S$ qui n'appartienne pas à \emph{l'ensemble de
ramification} de $\mathcal{R}$ : pour tout $S'$ dans $\mathrm{Et}(S)$ et tout
$M$ de $\mathcal{C}_{S'}$, le revêtement à opérateurs $\bigl(s(M),
\operatorname{Aut}(M)\bigr)$ de $S'$ doit être \emph{principal} aux points de
$S'$ au-dessus de l'image de $s$. Alors, pour tout objet $T$ de
$\mathrm{Rev}^{\mathcal{R}}(S)$, $r(T)$ est étale sur $S$ en $s$, et
\[
  F : T \longmapsto r(T)(s)
\]
est exact. Il est conservatif si et seulement si le site $\mathcal{C}$, muni de
la topologie image inverse de celle de $\mathrm{Et}(S)$, est \emph{connexe} ;
et si l'on suppose que $s(M)/\!\operatorname{Aut}(M) \to S'$ est un isomorphisme
pour tout $M$ --- « ce qui sera sans doute le cas dans tous les cas
intéressants » --- cela revient à dire que $S$ est connexe.

\subsection*{Le groupe}

Sous ces conditions, la théorie de Galois abstraite fournit un groupe profini
\[
  \pi_{1}^{\mathcal{R}}(S, s)
\]
et une équivalence entre $\mathrm{Rev}^{\mathcal{R}}(S)$ et la catégorie des
ensembles finis munis d'une action continue de ce groupe.

La connexité n'est pas un détail technique : c'est elle, et elle seule, qui
permet de remplacer un groupoïde par un groupe. Sans elle,
$\mathrm{Rev}^{\mathcal{R}}(S)$ n'est pas galoisienne mais seulement, comme la
page 30 le dira, \emph{multigaloisienne} --- un produit de catégories
galoisiennes indexé par les composantes connexes de l'objet final, donc un
\emph{groupoïde} profini et non un groupe.\footnote{Le mot apparaît page 30
sans définition, comme un terme déjà disponible : « Comme
$\mathrm{Rev}^{\mathcal{R}}(S)$ est en tous cas une catégorie
multigaloisienne \dots ». Sa définition axiomatique est aux pages 208 à 216 du
même carton, non datées ; ce sont les conditions A, B et C, et la
caractérisation (iii) y est précisément « $C$ équivalente à un produit
$\prod_{i} C_{i}$, les $C_{i}$ galoisiennes ». Le carton contient donc l'usage
du mot et son axiomatique, sans dire lequel des deux est venu d'abord, et cette
lecture ne le décide pas.}

\pagerange{25}{29}
\section*{Fonctorialité en $S$ et en $\mathcal{R}$ (pages 25 à 29)}

Soit $f : T \to S$ un morphisme de schémas, d'où un foncteur continu
$f^{\circ} : \mathrm{Et}(S) \to \mathrm{Et}(T)$. Donnons-nous une donnée
$\mathcal{R}_{S} = (\mathcal{C}_{S}, s_{S})$ sur $S$, une donnée
$\mathcal{R}_{T} = (\mathcal{C}_{T}, s_{T})$ sur $T$, un foncteur
$M \mapsto M \times_{S} T$ de $\mathcal{C}_{S}$ dans $\mathcal{C}_{T}$
transformant morphismes cartésiens en morphismes cartésiens et relevant
$f^{\circ}$, et enfin un morphisme entre les deux composés

\begin{tikzcd}[column sep=huge]
\mathcal{C}_{S} \arrow[r] \arrow[d, "s_{S}"'] & \mathcal{C}_{T} \arrow[d, "s_{T}"] \\
(\mathrm{Sch})_{/S} \arrow[r, "X \mapsto X \times_{S} T"'] & (\mathrm{Sch})_{/T}
\end{tikzcd}

soit
\[
  u(M) : s_{T}(M \times_{S} T) \longrightarrow s_{S}(M) \times_{S} T ,
\]
morphisme au-dessus de $f^{\circ}(S')$ pour chaque $M$ sur $S'$.

On en déduit un foncteur
$\mathrm{Rev}^{\mathcal{R}_{S}}(S) \to \mathrm{Rev}^{\mathcal{R}_{T}}(T)$ et un
homomorphisme de foncteurs composés

\begin{tikzcd}[column sep=huge]
\mathrm{Rev}^{\mathcal{R}_{S}}(S) \arrow[r] \arrow[d, "r"'] & \mathrm{Rev}^{\mathcal{R}_{T}}(T) \arrow[d, "r"] \\
\mathrm{Rev}(S) \arrow[r, "X \mapsto X \times_{S} T"'] & \mathrm{Rev}(T)
\end{tikzcd}

Cet homomorphisme est un isomorphisme dans certains cas : sous les conditions
de l'exemple 3, lorsque, pour tout sous-groupe $H$ de
$G = \operatorname{Aut}(M)$, le passage au quotient $s(M)/H$ commute au
changement de base $T \to S$ --- ce qui est vérifié si $T \to S$ est plat, ou
dans les cas kummériens a) pourvu que les images inverses des $D_{i}$ existent
et servent à définir la donnée sur $T$.

Le point de la construction est que \emph{le foncteur ainsi obtenu est exact}.
Il en résulte, par les sorites de la théorie de Galois abstraite, que si $t$
est un point géométrique de $T$ non ramifié pour $\mathcal{R}_{T}$ dont l'image
$s$ dans $S$ est non ramifiée pour $\mathcal{R}_{S}$, et si $u$ induit un
isomorphisme sur les fibres géométriques en $t$, on a un homomorphisme
canonique
\[
  \pi_{1}^{\mathcal{R}_{T}}(T, t) \longrightarrow
  \pi_{1}^{\mathcal{R}_{S}}(S, s) .
\]
Cela s'applique aux données kummériennes des trois types dès que $T \to S$ est
régulier, respectivement lisse, de sorte que les images inverses des $D_{i}$
vérifient encore les conditions de croisements normaux.

Sur la transitivité pour un composé $U \to T \to S$, le tapuscrit renonce ---
« on recule devant la tâche \dots , jetant provisoirement un voile pudique sur
ces abîmes » --- et cette lecture n'y supplée pas.\footnote{Le renoncement est
significatif et non anecdotique : c'est la 2-catégorialité de la situation qui
rend l'écriture pénible, et c'est exactement ce que les pages 196 à 198
finiront par assumer en posant que les domaines de ramification sur $S$ forment
une \emph{2-catégorie}, et que celles-ci forment pour $S$ variable une
2-catégorie fibrée.}

\subsection*{Deux cas particuliers}

Le premier est le changement de base \emph{étale}, avec $\mathcal{R}_{T}$
induite par $\mathcal{R}_{S}$ : le carré est déjà connu, puisqu'on dispose du
foncteur cartésien $(\ast)$ sur $\mathrm{Et}(S)$, et l'homomorphisme de
foncteurs composés y est un isomorphisme. Si de plus $T$ est fini sur $S$, avec
$S$ et $T$ connexes, et sous les conditions de l'exemple 3, l'homomorphisme sur
les groupes fondamentaux est un \emph{monomorphisme}, d'image le sous-groupe de
$\pi_{1}^{\mathcal{R}}(S,s)$ défini par le revêtement pointé $T$ de $S$.

Le second est le changement de base \emph{entier, radiciel et surjectif}, avec
$\mathcal{R}_{T}$ image inverse de $\mathcal{R}_{S}$. Alors le foncteur est
\emph{trivialement} une équivalence, par le résultat analogue pour les
revêtements étales de $X$ et $X \times_{S} T$ ; en particulier, sur un point
géométrique non ramifié de $T$, on obtient un isomorphisme des groupes
fondamentaux.

\subsection*{L'image inverse d'une donnée, et pourquoi elle est encombrante}

Grothendieck signale une difficulté qu'il ne résout pas. Si $\mathcal{R}$ est
définie par $(Z,G)$ comme à l'exemple 1, son image inverse par $T \to S$ est
définie par $(Z \times_{S} T,\, G)$ ; c'est immédiat. Mais si $\mathcal{R}$ est
la donnée kummérienne associée à un système $(D_{i})$, l'image inverse au sens
strict n'est \emph{pas} la donnée associée au système $(D_{i\,T})$, même à
équivalence de catégories fibrées près : elle lui est seulement
« essentiellement équivalente » au sens ci-dessus.

Il en tire deux propositions : ou bien introduire la notion d'image inverse
avec tout son poids, ou bien --- c'est le parti qu'il juge le plus raisonnable
--- donner des conditions axiomatiques sur le carré $(\ast)$, sous les
conditions de l'exemple 3, qui assurent \emph{moralement} que
$\mathcal{R}_{T}$ est essentiellement équivalente à l'image inverse de
$\mathcal{R}_{S}$, et qui s'appliquent directement aux cas kummériens. Il
ajoute qu'il serait raisonnable, « assez rapidement après les définitions
générales », de se borner aux données du type de l'exemple 3 : ce sont les
seules qu'on rencontrera jamais et les seules où l'on sache construire $r$.

\subsection*{Descente et passage à la limite}

Deux énoncés sont annoncés et non démontrés, avec l'usage qu'on veut en faire.

Un morphisme fidèlement plat et quasi-compact $T \to S$ est de \emph{descente
effective} pour la catégorie fibrée des revêtements de type $\mathcal{R}_{S}$
sur le diagramme $(S, T, T\times_{S}T, T\times_{S}T\times_{S}T)$. En
particulier la descente fidèlement plate vaut pour les revêtements kummériens,
ce qui est utilisé sans démonstration dans les notes de Murre pour la descente
à partir d'un anneau complété.

Il y a lieu ensuite d'énoncer un \emph{passage à la limite} pour un système
projectif $(X_{i})$ de schémas quasi-compacts et quasi-séparés à morphismes de
transition affines : cela ramène l'étude au voisinage d'un point $s$ à celle de
$\operatorname{Spec}(\mathcal{O}_{S,s})$. Ainsi, pour $S$ localement noethérien
normal muni d'une famille $(D_{i})$ de diviseurs réduits --- où la catégorie
des revêtements kummériens est une sous-catégorie de celle des revêtements ---
un revêtement $X$ de $S$ est kummérien si et seulement si le revêtement induit
sur $\operatorname{Spec}(\mathcal{O}_{S,s})$ l'est pour tout $s$, et l'on peut
y remplacer $\mathcal{O}_{S,s}$ par son hensélisé, par son hensélisé strict,
voire par son complété si les hypothèses restent stables par complétion --- par
exemple si $\mathcal{O}_{S,s}$ est excellent. Sans hypothèse de normalité ni de
réduction, les mêmes observations valent à condition de se borner aux
revêtements modérément ramifiés \emph{galoisiens}, ce qui est le sujet du
paragraphe suivant.

\pagerange{30}{31}
\section*{Revêtements galoisiens, et pourquoi la catégorie est indispensable (pages 30 et 31)}

\subsection*{La définition, complétée}

Un \emph{revêtement galoisien de type $\mathcal{R}$ de $S$, à groupe de Galois
$G$}, est un objet $X$ de $\mathrm{Rev}^{\mathcal{R}}(S)$ sur lequel le groupe
fini $G$ opère à droite de telle façon que, \emph{dans la catégorie}
$\mathrm{Rev}^{\mathcal{R}}(S)$, on obtienne un revêtement principal de l'objet
final : le morphisme $X \times G \to X \times X$ est un isomorphisme, et
$X \to e$ est couvrant.

Il faut y ajouter une condition que la page ne porte pas, sans quoi la
définition explicite qui suit ne coïncide pas avec celle-ci : \emph{l'action de
$G$ doit commuter aux actions des groupes d'automorphismes dont les modèles
sont déjà munis}, aussi bien qu'à la localisation étale sur $S$.\footnote{C'est
la question de Murre, page 2 : « It seems to me that one must require that the
actions of $G$ and $\mathcal{G}$ commute; is this correct? » La réponse, page
36, est oui : « the action of $G$ commutes to étale localization on $S$
\emph{and} to the action of the groups of automorphisms given on the various
$Z$'s ». Le paragraphe 4 lui-même ne la porte pas ; elle est rétablie ici, à sa
place logique, parce que sans elle l'explicitation qui suit --- et le critère
de pleine fidélité de la page 31 --- ne veut rien dire. Voir aussi la lecture
du lot 1.}

Sous cette condition, la notion s'explicite de trois façons équivalentes :
comme ci-dessus ; en termes des groupes fondamentaux, puisque
$\mathrm{Rev}^{\mathcal{R}}(S)$ est en tous cas une catégorie
multigaloisienne ; et en revenant à la définition, en exigeant que, pour tout
$M \in \mathcal{C}_{S'}$, le revêtement $X(M)$ de $s(M)$ soit \emph{principal
de groupe $G$}.

Par abus de langage, on appellera aussi revêtement galoisien de type
$\mathcal{R}$ le revêtement à opérateurs $\bigl(r(X), G\bigr)$. L'abus est sans
danger lorsque le foncteur de la catégorie des revêtements galoisiens de type
$\mathcal{R}$ --- morphismes restreints aux \emph{isomorphismes} --- dans celle
des revêtements à opérateurs de $S$ est pleinement fidèle : alors, pour un
revêtement à opérateurs $(T,G)$ donné, il existe au plus un $(X,G)$ galoisien
de type $\mathcal{R}$ dont il provienne, à isomorphisme canonique près.

\subsection*{La condition nécessaire et suffisante}

Voici la condition, d'abord dans le cas de l'exemple 1. Pour tout localisé
strict $S'$ d'un point de $S$, soit $Z'_{0}$ une composante connexe de
$Z' = Z \times_{S} S'$ et $H$ son stabilisateur dans $G$. On demande : pour
deux sous-groupes distingués $H_{1}, H_{2}$ de $H$ et tout isomorphisme
$u : Z'_{0}/H_{1} \to Z'_{0}/H_{2}$ compatible avec un isomorphisme
$\varphi : H/H_{1} \to H/H_{2}$, il existe $g \in H$ tel que
$\operatorname{int}(g)$ applique $H_{1}$ sur $H_{2}$, induise $\varphi$ par
passage au quotient, et tel que $u$ soit induit par $g_{Z'_{0}}$.

Dans le cas général de l'exemple 3, on exige cela pour tous les couples
$\bigl(s(M), \operatorname{Aut}(M)\bigr)$ sur les divers $S'$ de
$\mathrm{Et}(S)$.

Cette condition est satisfaite dans le cas kummérien --- Grothendieck renvoie
aux énoncés 1.10 et 1.15 des notes de Murre --- et la démonstration du critère
est esquissée page 37 : la question est locale pour la topologie étale, donc on
se ramène au cas strictement local, où la classification est explicite (voir
plus bas).

\subsection*{Le passage en anglais}

C'est ici que le tapuscrit change de langue, en pleine page, parce que
Grothendieck est en train de répondre à Murre. Le passage est le centre du
dossier et il tient en quatre affirmations.

\begin{enumerate}
  \item Ce phénomène explique pourquoi il a été si réticent --- à tort ---
    à définir les revêtements modérément ramifiés hors du cas galoisien :
    \emph{dans le cas non galoisien, ce ne sera pas un revêtement satisfaisant
    des conditions supplémentaires, mais un revêtement muni d'une structure
    supplémentaire} ; alors que dans le cas galoisien, la structure est
    uniquement déterminée si elle existe, et l'on peut ne pas la voir.
  \item Tous les morphismes au sens naïf entre revêtements modérés, même
    galoisiens, ne doivent pas être considérés : seuls ceux qui viennent de
    morphismes de $\mathrm{Rev}^{\mathcal{R}}(S)$. Le contre-exemple à 1.16
    (lot 1) montre qu'il faut y prendre garde.
  \item La notion naturelle de \emph{produit} de deux revêtements modérés,
    galoisiens ou non, n'est pas le produit fibré naïf au-dessus de $S$ ---
    ce qui, ajoute-t-il, explique la complication de la démonstration d'un des
    lemmes de Murre.
  \item Donc : « This shows rather clearly that even in the Galois case, the
    consideration of the category $\mathrm{Rev}^{\mathcal{R}}(S)$ is
    undispensable for a good comprehension of what is really going on. »
\end{enumerate}

\pagerange{32}{33}
\section*{Le troisième état de l'exemple 3 : cinq conditions sur $\mathcal{R}$ (pages 32 et 33)}

Deux feuillets séparés reprennent l'exemple 3 et l'ouvrent par un jugement :
« La présentation de cet exemple n'est ni élégante ni naturelle, la donnée de
$s$ et de \dots\ étant un élément de structure artificiel. » Suit la
présentation intrinsèque, en cinq conditions sur $\mathcal{R}$ seule.

\begin{enumerate}
  \item[a)] Pour tout objet $M$ d'un $\mathcal{C}_{S'}$, le faisceau étale
    associé au préfaisceau des automorphismes de $M$ sur $\mathrm{Et}(S')$ est
    représentable par un schéma en groupes $G(M)$ fini et étale sur $S'$,
    déterminé à isomorphisme unique près. Il opère sur le schéma relatif
    $s(M)$.
  \item[b)] Pour tout $S'$, tout morphisme de $\mathcal{C}_{S'}$ est un
    épimorphisme et tout endomorphisme un automorphisme ; et pour tout
    $u : M' \to M$ et tout automorphisme $g'$ de $M'$, il existe un
    automorphisme $g$ de $M$ tel que $gu = ug'$. Cet automorphisme est alors
    unique, et $u$ détermine un unique homomorphisme de schémas en groupes
    $G(M') \to G(M)$ le rendant équivariant.
  \item[c)] Cet homomorphisme $G(M') \to G(M)$ est surjectif, et si $H$ en est
    le noyau, $s(M') \to s(M)$ induit un isomorphisme $s(M')/H \simeq s(M)$ ---
    les $s(M)$ étant supposés relativement affines, de sorte que le quotient a
    un sens.
  \item[d)] Deux homomorphismes de $M'$ dans $M$ se déduisent l'un de l'autre,
    localement pour la topologie étale, par composition avec un automorphisme
    de $M$.
  \item[e)] Étant donnés $M$, $M'$ dans $\mathcal{C}_{S'}$, il existe une
    famille couvrante $S'_{i} \to S'$ dans $\mathrm{Et}(S)$ telle que les
    images de $M$ et $M'$ sur $S'_{i}$ soient dominées par un même objet
    $M''_{i}$ de $\mathcal{C}_{S'_{i}}$.
\end{enumerate}

Trois choses sont à voir. La première est le mouvement : le groupe
d'automorphismes, qui n'était qu'un groupe abstrait à la première mouture et un
groupe fini à la seconde, est ici un \emph{schéma en groupes}, obtenu par
faisceautisation du préfaisceau des automorphismes --- donc extrait de
$\mathcal{C}$ et non ajouté à elle. La condition c) fait alors de $s(M') \to
s(M)$ un torseur sous $H$, ce que la page ébauche avant de biffer.

La deuxième est ce que cela permet : « on peut, dans la description des données
de ramification kummériennes, se dispenser d'introduire aussi des isomorphismes
$u : G_{S'} \to \mu_{\underline{n},\,S'}$ ». Le quadruple
$(\underline{a}, \underline{n}, G, u)$ du lot 1 se réduit au couple
$(\underline{a}, \underline{n})$, et c'est la réponse à la question de Murre.

La troisième est la conditions d) et e) prises ensemble : $\mathcal{C}_{S'}$
est localement un groupoïde connexe. C'est la structure que les pages 198 à 200
appelleront une \emph{gerbe}, et sous cette forme la donnée kummérienne est la
gerbe des racines $\underline{n}$-ièmes des équations des $D_{i}$.

\pagerange{34}{38}
\section*{La lettre d'accompagnement : le cas strictement local (pages 34 à 38)}

La lettre qui accompagne le tapuscrit est une explication, non un complément :
elle reprend l'exemple 1 dans le langage des corps de fonctions (lu au lot 1),
puis fait le calcul strictement local, qui est ce que l'esquisse omet.

Plaçons-nous sur $S$ strictement local. Pour des revêtements définis sur $S$
tout entier, la situation générale de l'exemple 3 se ramène à celle de
l'exemple 2 ; et, quitte à remplacer les $Z_{i}$ par des composantes connexes,
on peut les supposer connexes. Pour $Z_{i}$ fixé, les revêtements étales de
$Z_{i}$ sont triviaux, puisque $Z_{i}$ est lui aussi strictement local. Donc :
\[
  \bigl\{\text{revêtements étales de } (Z_{i}, G_{i})\bigr\}
  \;\simeq\;
  \bigl\{\text{ensembles finis } E \text{ munis d'une action de } G_{i}\bigr\} ,
\]
et la réalisation géométrique sur $S$ de l'objet associé à $E$ est
$Z_{i} \times^{G_{i}} E$.

Cela donne, pour $S$ quelconque et localement pour la topologie étale, une
description concrète de ce que signifie, pour un revêtement $S'$ de $S$, d'être
associé à un modèle $M$ : \emph{cela signifie qu'il est localement décrit par
une représentation d'un groupe d'inertie sur un ensemble fini} ; et la
structure supplémentaire consiste à se donner de telles représentations locales
qui se recollent. La description a un sens dès que $r$ est fidèle, sans qu'on
ait à le supposer pleinement fidèle.

Pour les revêtements galoisiens de groupe $H$, le calcul est le suivant, et il
est joli. Dans le cas strictement local, pour $Z_{i}$ fixé, la catégorie de
ces revêtements est équivalente à celle des ensembles finis $E$ munis d'une
action à gauche de $G_{i}$ et d'une action à droite de $H$ qui commutent, $E$
étant un torseur sous $H$. À isomorphisme près, ils sont classés par
\[
  \operatorname{Hom}(G_{i}, H) \big/ H ,
\]
homomorphismes de groupes modulo automorphisme intérieur.\footnote{Le calcul
est immédiat et vaut d'être refait : $E$ étant un $H$-torseur, le choix d'un
point l'identifie à $H$ opérant sur lui-même à droite ; les automorphismes de
$H$ comme $H$-ensemble à droite sont les translations à gauche ; une action à
gauche de $G_{i}$ qui commute aux translations à droite est donc donnée par un
homomorphisme $G_{i} \to H$, et le changement du point choisi le conjugue.}

C'est de cette description que sort le critère de pleine fidélité du foncteur
« réalisation géométrique » sur les revêtements galoisiens de type
$\mathcal{R}$ : la question est locale pour la topologie étale, ce qui ramène
au cas strictement local, où l'on vient de tout calculer.

La lettre porte enfin, en passant, une remarque de vocabulaire qui n'est pas
anodine : « By the way, one should rather say \emph{category of ramification
data} ». Le carton finira par lui donner raison.

\pagerange{39}{40}
\section*{Un dossier s'ouvre : « Finitude (Pbs et conjectures) » (pages 39 et 40)}

Les deux dernières pages du lot ouvrent une chemise nouvelle, de la main de
Grothendieck, et sans rapport apparent avec ce qui précède --- sauf qu'elle
demande, du groupe fondamental, ce que le tapuscrit vient de construire.

Trois énoncés, sur un corps $k$ algébriquement clos.

\begin{enumerate}
  \item[a)] $X$ propre sur $k$ $\Longrightarrow$ $\pi_{1}(X)$ topologiquement
    de type fini.
  \item[b)] $X$ de type fini sur $k$ $\Longrightarrow$ $\pi'_{1}(X)$ de
    \emph{présentation finie}, où $\pi'_{1}$ désigne le plus grand quotient
    premier à la caractéristique.
  \item[c)] $X$ de type fini sur $k$ $\Longrightarrow$ $\pi^{t}_{1}(X)$
    topologiquement de type fini, où $\pi^{t}_{1}$ classifie les revêtements
    modérément ramifiés le long du bord d'une compactification.
\end{enumerate}

Et une note en marge : « c) plus général que a) ». En effet, pour $X$ propre il
n'y a pas de bord, donc $\pi^{t}_{1}(X) = \pi_{1}(X)$, et c) redonne
a).\footnote{Le sens précis de $\pi_{1}^{t}$ est indiqué en marge de la page
d'une écriture très abîmée ; la transcription y laisse plusieurs mots
illisibles. La lecture retenue ici est celle que le reste du carton impose ---
notamment les pages 112 à 115, où $\pi_{1}^{t}(U)$ est défini comme le quotient
de $\pi_{1}(U)$ classifiant les revêtements modérément ramifiés le long des
composantes de codimension 1 du bord.}

L'état des démonstrations est noté en regard, et il faut le lire comme un
tableau de bord. L'énoncé a) est \emph{démontré} : réduction au cas normal,
Bertini pour se ramener à une courbe, puis théorie de la spécialisation du
groupe fondamental et méthode transcendante. Pour b), « le cas type fini est
prouvé » : réduction au cas normal par descente, puis au cas où les
singularités sont « jolies au sens Artin », suite exacte des $\pi_{1}$,
c'est-à-dire théorème de spécialisation, puis passage à une courbe normale.
Reste ce que la résolution des singularités permettrait de faire, et le cas
d'une surface projective avec un diviseur à croisements normaux.

Ces trois questions occupent les lots 3 et 4, où l'on trouvera le théorème de
finitude relatif, le groupe fondamental birationnel, et l'énoncé --- démontré
là --- que $\pi'_{1}(X \times Y) = \pi'_{1}(X) \times \pi'_{1}(Y)$.

\end{document}
